\documentclass[11pt]{beamer}
\usepackage[utf8]{inputenc}
\usepackage{amsmath, amssymb, graphicx}
\usepackage{bm}

% Presentation Theme
\usetheme{Madrid}
\usecolortheme{default}

\title[Gravity Assists \& Flybys]{Astrodynamics 2026-1}
\subtitle{Planetary Flybys \& Gravity Assist Maneuvers}
\author{Prof. Juan}
\institute{Universidad de Antioquia - Aerospace Engineering}
\date{Spring 2026}

\begin{document}

% --- Title Slide ---
\begin{frame}
    \titlepage
\end{frame}

% ==========================================
% SECTION 1: THE FLYBY CONCEPT
% ==========================================
\section{The Flyby Concept}

% --- Slide 1 ---
\begin{frame}{What is a Planetary Flyby?}
    A spacecraft that enters a planet's sphere of influence (SOI) and does not impact the planet or go into orbit around it will continue in its hyperbolic trajectory through periapsis and exit the sphere of influence.
    
    \vspace{0.3cm}
    \textbf{The Patched Conic View:}
    \begin{itemize}
        \item \textbf{Inbound:} The spacecraft arrives at the SOI with a hyperbolic excess velocity, $\bm{v}_{\infty_1}$.
        \item \textbf{Outbound:} It exits the SOI with a hyperbolic excess velocity, $\bm{v}_{\infty_2}$.
        \item \textbf{Energy Conservation:} Relative to the planet, energy is conserved. Therefore, the spacecraft exits the SOI with the exact same relative speed as it entered: $|\bm{v}_{\infty_1}| = |\bm{v}_{\infty_2}| = v_\infty$.
    \end{itemize}
\end{frame}

% --- Slide 2 ---
\begin{frame}{The Turn Angle ($\delta$)}
    While the \textit{speed} relative to the planet does not change, the \textit{direction} of the velocity vector is rotated by the planet's gravity.
    
    \vspace{0.3cm}
    This rotation is defined by the \textbf{Turn Angle} ($\delta$), which represents the angle between the incoming and outgoing asymptotes of the hyperbola:
    \begin{equation}
        \delta = 2 \sin^{-1}\left(\frac{1}{e}\right)
    \end{equation}
    
    Where the eccentricity ($e$) of the flyby hyperbola is dictated by how close the spacecraft passes to the planet (the periapsis radius, $r_p$):
    \begin{equation}
        e = 1 + \frac{r_p v_\infty^2}{\mu}
    \end{equation}
   
\end{frame}

% ==========================================
% SECTION 2: VECTOR MECHANICS
% ==========================================
\section{Vector Mechanics of a Gravity Assist}

% --- Slide 3 ---
\begin{frame}{Heliocentric Velocity Change}
    The magic of a gravity assist happens when we shift our frame of reference back to the Sun. 
    
    \vspace{0.3cm}
    Let $\bm{V}$ be the heliocentric velocity of the planet. The heliocentric velocity of the spacecraft at the inbound and outbound SOI crossings are:
    \begin{align}
        \bm{V}_1^{(v)} &= \bm{V} + \bm{v}_{\infty_1} \\
        \bm{V}_2^{(v)} &= \bm{V} + \bm{v}_{\infty_2}
    \end{align}
   
    
    The absolute change in the spacecraft's heliocentric velocity ($\Delta \bm{V}^{(v)}$) is simply the vector difference of the asymptotes:
    \begin{equation}
        \Delta \bm{V}^{(v)} = \bm{V}_2^{(v)} - \bm{V}_1^{(v)} = \bm{v}_{\infty_2} - \bm{v}_{\infty_1} = \Delta \bm{v}_\infty
    \end{equation}
   
\end{frame}

% --- Slide 4 ---
\begin{frame}{The Trailing-Side Flyby (Speed Up)}
    \textbf{Trajectory Geometry:} The spacecraft crosses the orbit of the planet \textit{behind} the planet (trailing side).
    
    \vspace{0.2cm}
    \begin{itemize}
        \item The planet's gravity pulls the spacecraft "forward" as it passes behind.
        \item The turn angle ($\delta$) rotates $\bm{v}_{\infty_2}$ so it aligns more closely with the planet's velocity vector $\bm{V}$.
        \item \textbf{Result:} The component of $\Delta \bm{V}^{(v)}$ in the direction of the planet's velocity is positive, meaning a trailing-side flyby \textbf{increases} the spacecraft's heliocentric speed.
    \end{itemize}
\end{frame}

% --- Slide 5 ---
\begin{frame}{The Leading-Side Flyby (Slow Down)}
    \textbf{Trajectory Geometry:} The spacecraft crosses the orbit of the planet \textit{in front} of the planet (leading side).
    
    \vspace{0.2cm}
    \begin{itemize}
        \item The planet's gravity slows the spacecraft down as it passes in front.
        \item The turn angle ($\delta$) rotates $\bm{v}_{\infty_2}$ away from the planet's velocity vector $\bm{V}$.
        \item \textbf{Result:} The component of $\Delta \bm{V}^{(v)}$ in the direction of the planet's velocity is negative, meaning a leading-side flyby \textbf{decreases} the spacecraft's heliocentric speed.
    \end{itemize}
\end{frame}

% ==========================================
% SECTION 3: REAL-WORLD APPLICATIONS
% ==========================================
\section{Real-World Applications}

% --- Slide 6 ---
\begin{frame}{Historical Gravity Assists}
    Gravity assist maneuvers are used to add momentum to a spacecraft over and above that available from a spacecraft's onboard propulsion system, granting access to regions of the solar system that would otherwise be inaccessible.
    
    \vspace{0.2cm}
    \textbf{Famous Examples:}
    \begin{itemize}
        \item \textbf{Voyager 1 \& 2 (1977):} Used Jupiter and Saturn flybys to achieve solar system escape velocity.
        \item \textbf{Galileo (1989):} Too heavy to fly directly to Jupiter, it used a Venus-Earth-Earth Gravity Assist (VEEGA) sequence.
        \item \textbf{Cassini (1997):} Arrived at Saturn in 2004 using four flybys in a Venus-Venus-Earth-Jupiter (VVEJGA) sequence.
    \end{itemize}
\end{frame}

% ==========================================
% SECTION 4: ANALYTICAL EXAMPLE
% ==========================================
\section{Analytical Example: Venus Flyby}

% --- Slide 7 ---
\begin{frame}{Example: Approaching Venus}
    Let's analyze a spacecraft arriving at Venus for a gravity assist to drop its perihelion closer to the sun. We are targeting a \textbf{300 km altitude} flyby on the \textbf{dark side} (leading side) of the planet.
    
    \vspace{0.2cm}
    \textbf{Given Parameters:}
    \begin{itemize}
        \item $\mu_{Venus} = 324,900 \text{ km}^3/\text{s}^2$
        \item $r_{Venus} = 6052 \text{ km}$
        \item Flyby Periapsis ($r_p$) = $6052 + 300 = \mathbf{6352 \text{ km}}$
        \item Spacecraft Inbound Heliocentric Speed ($V_1^{(v)}$) = $\mathbf{37.62 \text{ km/s}}$
        \item Venus Heliocentric Speed ($V$) = $\mathbf{35.02 \text{ km/s}}$
    \end{itemize}
\end{frame}

% --- Slide 8 ---
\begin{frame}{Example: The Encounter Hyperbola}
    First, we determine the geometry of the hyperbolic encounter relative to Venus. Vector addition tells us the incoming hyperbolic excess speed ($v_\infty$) is $\mathbf{3.733 \text{ km/s}}$.
    
    \vspace{0.2cm}
    \textbf{1. Calculate Eccentricity ($e$):}
    \begin{equation}
        e = 1 + \frac{r_p v_\infty^2}{\mu_{Venus}} = 1 + \frac{6352 (3.733)^2}{324,900} = \mathbf{1.272}
    \end{equation}
   
    
    \textbf{2. Calculate Turn Angle ($\delta$):}
    \begin{equation}
        \delta = 2 \sin^{-1}\left(\frac{1}{e}\right) = 2 \sin^{-1}\left(\frac{1}{1.272}\right) = \mathbf{103.6^\circ}
    \end{equation}
   
\end{frame}

% --- Slide 9 ---
\begin{frame}{Example: The Heliocentric Result}
    By passing on the dark side (leading side), the spacecraft's outgoing asymptote is rotated $103.6^\circ$ \textit{away} from the planet's velocity vector. 
    
    \vspace{0.2cm}
    When we translate this back to the Sun-centered frame using vector addition ($V_2^{(v)} = V + v_{\infty_2}$), we find:
    
    \begin{itemize}
        \item \textbf{Inbound Heliocentric Speed:} $37.62 \text{ km/s}$
        \item \textbf{Outbound Heliocentric Speed:} $\mathbf{31.78 \text{ km/s}}$
    \end{itemize}
    
    \vspace{0.2cm}
    \textbf{Conclusion:} 
    Without firing a single engine, the spacecraft successfully scrubbed $\mathbf{5.84 \text{ km/s}}$ off its heliocentric speed, dropping its new orbit deeper into the inner solar system!
\end{frame}

% --- Slide: Insertion vs. Assist ---
\begin{frame}{Comparing the Patched Conics}
    Both Orbit Insertions and Gravity Assists rely on hyperbolic approach geometry, but they exploit different parameters.
    
    \vspace{0.3cm}
    \textbf{Orbit Insertion (Planetary Capture)}
    \begin{itemize}
        \item \textbf{Goal:} Transition from hyperbola ($e > 1$) to ellipse ($e < 1$).
        \item \textbf{Key Variable:} The \textbf{$\beta$ angle}.
        \item \textbf{Reason:} $\beta$ defines the physical location of periapsis relative to the incoming asymptote. This dictates the exact timing and location of the critical insertion burn. The $\delta$ angle is ignored because the outbound asymptote is never reached.
    \end{itemize}

    \vspace{0.3cm}
    \textbf{Gravity Assist (Planetary Flyby)}
    \begin{itemize}
        \item \textbf{Goal:} Rotate the $v_\infty$ vector to change heliocentric energy.
        \item \textbf{Key Variable:} The \textbf{$\delta$ angle (Turn Angle)}.
        \item \textbf{Reason:} $\delta$ defines the angular deflection of the trajectory. The larger the $\delta$, the more efficiently we can align our outbound vector with the planet's heliocentric velocity to steal momentum.
    \end{itemize}
    
    \vspace{0.2cm}
    \textit{Mathematical Relationship:} $\delta = 180^\circ - 2\beta$
\end{frame}

\end{document}